\documentclass[11pt,a4paper]{article}

\usepackage[margin=1in]{geometry}
\usepackage{xcolor}
\usepackage{parskip}
\usepackage{enumitem}
\usepackage{hyperref}
\usepackage{framed}
\usepackage{soul}
\usepackage{amsmath,amssymb}
\usepackage{siunitx}
\usepackage[authoryear,round]{natbib}

\hypersetup{
  colorlinks=true,
  linkcolor=blue!60!black,
  urlcolor=blue!60!black,
}

% ---------------------------------------------------------------
% Response-letter environments and macros
% ---------------------------------------------------------------

% Reviewer's comment: left bar, italic, grey
\definecolor{reviewerbar}{gray}{0.55}
\definecolor{reviewerbg}{gray}{0.96}

\newenvironment{rcomment}
  {\begin{leftbar}\color{black!75}\itshape}
  {\end{leftbar}\vspace{0.3em}}

% Author's response: normal style, marked with a bold tag
\newcommand{\response}{\par\noindent\textbf{Response:}\ \ignorespaces}

% Manuscript change pointer
\newcommand{\changed}[1]{\par\noindent\textbf{Manuscript changes:}\ #1\par}

% Cross-reference helpers
\newcommand{\manref}[1]{(#1)}                  % e.g., \manref{p.~4, l.~112}
\newcommand{\secref}[1]{Section~#1}
\newcommand{\figref}[1]{Fig.~#1}
\newcommand{\eqnref}[1]{Eq.~#1}

% Commonly needed symbols in this paper
\newcommand{\Kzero}{K_0}
\newcommand{\Zm}{Z_m}
\newcommand{\Zmz}{Z_{m0}}
\newcommand{\Vs}{V_s}
\newcommand{\pmean}{p^{\prime}}
\newcommand{\ru}{r_u}

% ---------------------------------------------------------------

\title{Response to Reviewers\\
  \large Manuscript COMGE-D-26-01109}
\author{Yusong Han, Ryosuke Uzuoka, Kyohei Ueda}
\date{\today}

\begin{document}

\maketitle

\section*{Cover letter}

Dear Editor,

Thank you for the opportunity to revise our manuscript
``\emph{[paper title]}'' (COMGE-D-26-01109) submitted to
\emph{Computers and Geotechnics}. We sincerely thank both reviewers for their
thoughtful and constructive comments, which have helped us to substantially
improve the manuscript.

We have carefully addressed every comment. A point-by-point response is
provided below. In the revised manuscript, added or modified text is marked
using \texttt{latexdiff} against the submitted version (also provided).

%% -------- Summary of major changes --------

\section*{Summary of major revisions}

\begin{itemize}[leftmargin=*]
  \item \textbf{Clarified the relative-density definition} as a
    behavior-matching label calibrated against \citet{...} and
    \citet{...} data, and added a new CSR=0.40 cyclic validation point at
    $D_r \approx 75\%$.
  \item \textbf{Expanded the Introduction} with a review of prior DEM-HCA
    models and the boundary-condition taxonomy (rigid wall, stacked wall,
    flexible membrane).
  \item \textbf{Added a consolidated case table} listing all
    ($D_r$, $\Kzero$, CSR) combinations and the resulting $N_L$.
  \item \textbf{Added a quantitative quasi-static check} through the
    inertial number.
  \item \textbf{Added contact-force anisotropy analysis} to complement the
    existing contact-density analysis.
  \item \textbf{Quantified the differences} in Fig.~15 via ... .
  \item Minor clarifications and additional references throughout.
\end{itemize}

All line and figure numbers below refer to the revised manuscript unless
noted otherwise.

\clearpage

%% ===============================================================
%% Reviewer 1
%% ===============================================================

\section*{Response to Reviewer 1}

We thank Reviewer~1 for the detailed and constructive feedback. Our
responses to each point are given below.

%% ---------------- R1.1 ----------------
\subsection*{R1.1 --- Particle size scaling adequacy}

\begin{rcomment}
Lines 112--114: Particle size may influence macroscopic mechanical responses
such as dilatancy, peak friction angle, and liquefaction resistance. Please
clarify the adequacy of this scaling factor (e.g., whether the specimen
size-to-particle size ratio meets the requirement) or provide supporting
references.
\end{rcomment}

\response
We thank the reviewer for raising this important point. The specimen is an
annular hollow cylinder with inner diameter \SI{60}{mm}, outer diameter
\SI{100}{mm}, height \SI{100}{mm}, and a particle-size distribution of
\SIrange{1.0}{3.0}{mm} ($d_{\max}=\SI{3}{mm}$). The radial thickness
$R-r=\SI{20}{mm}$ is inherently the smallest coherent dimension of any HCA
specimen; we acknowledge this openly as a structural feature of the test
geometry rather than a modelling oversight. The corresponding specimen-to-%
$d_{\max}$ ratios are 6.7 radially, 33 axially, and 84 in the mid-radius
circumferential direction. We address the adequacy of the radial ratio
through the following independent considerations.

\begin{enumerate}[leftmargin=*,label=(\arabic*)]
  \item \textbf{Codified laboratory guideline.} The most geometrically
    analogous standardised guideline for cohesive-soil triaxial testing is
    ASTM~D4767, which stipulates that the largest particle size shall be
    smaller than one-sixth of the specimen diameter. The present radial
    ratio of 6.7 satisfies this 6:1 threshold; the axial and circumferential
    ratios satisfy it by a wide margin.

  \item \textbf{Consistency with DEM-HCA peer literature.} Among prior DEM-%
    HCA studies employing rigid-wall boundaries, \citet{Shi2021}, using a
    bonded-particle model of sandstone under HCA torsional loading, adopted
    a radial/$d_{\max}$ ratio of 7.5 --- directly comparable to the present
    6.7. Separately, the 10$\times$ up-scaling of the Toyoura grading adopted
    here follows the same computational-efficiency rationale as in the DEM-%
    HCA Toyoura-sand study of \citet{Song2024}. The present choices ---
    both in size ratio and in scaling strategy --- are therefore consistent
    with peer-reviewed DEM-HCA practice.

  \item \textbf{Liquefaction is an end-state problem, and our comparison
    is differential.} Liquefaction is defined at a finite strain level
    ($\gamma_{SA}=2.5\%$ in the \citealt{Ishihara1985} protocol, which we
    follow) and is governed by end-state fabric collapse rather than small-%
    strain elastic response. Classical critical-state soil mechanics holds
    that ultimate shear strength is a material property that is independent
    of specimen size. Moreover, the principal interest of the present work
    is the \emph{variation} of $N_L$ across \Kzero{} states rather than its
    absolute value: any size-induced bias applies uniformly across the
    compared \Kzero{} states and therefore does not distort the relative
    trends that form the paper's conclusions.

  \item \textbf{HCA-specific geometric argument.} The primary shear strain
    $\gamma_{z\theta}$ developed during cyclic torsional loading lies in the
    tangential--axial plane. The smallest dimension (radial) is therefore
    \emph{perpendicular} to the primary deformation mode, unlike direct-%
    shear tests in which loading acts across the thinnest dimension. Wall-%
    induced inhomogeneity in the radial direction affects
    $\sigma'_r$ measurement, which we address through annulus averaging,
    but does not directly distort the primary $\gamma_{z\theta}$ response.
    Representative-volume-element guidelines developed for cubic or solid-%
    cylindrical specimens under axial loading do not apply cleanly to the
    HCA configuration.

  \item \textbf{Empirical corroboration.} The calibrated parameter set
    simultaneously reproduces three independent experimental behaviours of
    Toyoura sand: the monotonic effective stress path and stress--strain
    response of \citet{Nakata1998}, the associated dilatancy, and the
    cyclic liquefaction-resistance curve of \citet{Ishihara1985} across
    five CSR levels. This triple behavioural match provides empirical
    evidence that the adopted particle sizing is adequate for the stress
    regime and loading mode investigated.

  \item \textbf{Computational trade-off.} A further reduction of the
    particle size (e.g., halving $d_{50}$ to meet the more stringent
    representative-volume-element ratios developed for axial-compression
    specimens) would increase the computational cost by approximately 16$\times$
    --- an 8-fold increase in particle number in three dimensions combined
    with an approximately 2-fold reduction in the critical timestep ---
    with limited incremental scientific return for a comparative study of
    \Kzero{} effects. We note this refinement as a direction for future
    work rather than pursue it here.
\end{enumerate}

\changed{In the revised manuscript we have (i) expanded the DEM-HCA
literature review in the Introduction (Section~\ref{sec:introduction}, last
paragraph) to introduce \citet{LiGuoZhang2014} (the reference specifically
highlighted by the reviewers), \citet{Shi2021}, and \citet{Song2024}
alongside the previously cited \citet{LiZhangGutierrez2015} and
\citet{Liu2021}, and (ii) added a brief note in
Section~\ref{subsec:dem_hca} citing \citet{Song2024} as precedent for the
10$\times$ scaling of the Toyoura-sand grading.}

%% ---------------- R1.2 ----------------
\subsection*{R1.2 --- Void ratio range and ``dense''/``medium-dense'' labels}

\begin{rcomment}
For the case where the authors have defined relative density states of
$D_r \approx 90\%$ and $D_r \approx 75\%$ based on void ratios of 0.736 and
0.742, respectively, the implied range of maximum and minimum void ratios is
approximately 0.732 to 0.772. As a result, I do not think that a void ratio
range of merely 0.04 and the terms ``dense'' and ``medium-dense'' as used in
this manuscript correspond to the physical states observed in the Toyoura
sand that this model purports to represent.
\end{rcomment}

\response
\textit{[TODO: reframe $D_r$ as a behavior-matching label; cite Nakata
monotonic + Ishihara cyclic triple match; report new CSR=0.40 $D_r\approx
75\%$ validation point; acknowledge absence of Nakata monotonic data at
$D_r \approx 75\%$ as a limitation.]}

\changed{\textit{[TODO: \secref{...}, revised \figref{7}, new paragraph on
limitations.]}}

%% ---------------- R1.3 ----------------
\subsection*{R1.3 --- Missing formulas for CSR and $\Kzero$}

\begin{rcomment}
Several variable symbols are used without a clear computational formula.
Such as CSR, $\Kzero$.
\end{rcomment}

\response
\textit{[TODO: add CSR $= \tau_{cyc}/\sigma'_{z0}$ and
$\Kzero = \sigma'_r/\sigma'_z$ at first use.]}

\changed{\textit{[TODO: pointer to \secref{2.3}, lines X--Y.]}}

%% ---------------- R1.4 ----------------
\subsection*{R1.4 --- Consolidated case table}

\begin{rcomment}
The description of the simulation program is summarized in brief statements
without explicitly listing the complete set of parameters. It is recommended
that all simulation cases (including combinations of different densities,
$\Kzero$ values, and CSR levels) be presented collectively in a table
format.
\end{rcomment}

\response
\textit{[TODO: confirm new table with $D_r$, $\Kzero$, CSR, $N_L$.]}

\changed{\textit{[TODO: new Table~\ref{...}.]}}

%% ---------------- R1.5 ----------------
\subsection*{R1.5 --- Loading rate / quasi-static condition}

\begin{rcomment}
The loading rate is not specified in the text. In DEM simulations, an
excessively high loading rate may introduce inertial effects, leading to
excessive particle accelerations and distorted contact forces. Please
clarify whether the selected loading rate satisfies quasi-static conditions.
\end{rcomment}

\response
\textit{[TODO: report inertial number $I = \dot{\gamma}\,d\,\sqrt{\rho/p}$;
show $I \ll 10^{-3}$; cite da Cruz / Roux.]}

\changed{\textit{[TODO: \secref{2.3}, new equation / paragraph.]}}

%% ---------------- R1.6 ----------------
\subsection*{R1.6 --- Rigid wall vs.\ flexible membrane}

\begin{rcomment}
Rigid walls are used to model the inner and outer boundaries of the HCA
specimen, whereas flexible latex membranes are typically used in laboratory
experiments. Rigid boundaries may restrict local deformation and shear band
development, potentially leading to overestimation of stiffness and
liquefaction resistance. Please discuss the limitations of adopting rigid
boundaries and justify their use with relevant references.
\end{rcomment}

\response
\textit{[TODO: acknowledge limitation; justify via (i) HCA chambers are
rigid-walled with the membrane as a pressure interface, (ii) present work
focuses on average fabric-level comparison rather than localized shear
bands, (iii) flexible-membrane DEM-HCA (Song 2024) as complementary / future
work.]}

\changed{\textit{[TODO: \secref{2.1} end paragraph or discussion section.]}}

%% ---------------- R1.7 ----------------
\subsection*{R1.7 --- Citation for CSL in Fig.~8(c)}

\begin{rcomment}
The CSL in Fig.~8(c) is of unknown origin; whether it comes from your own
data or from someone else's literature, a clear citation is required.
\end{rcomment}

\response
\textit{[TODO: verify source; update caption with explicit citation.]}

\changed{\textit{[TODO: \figref{8}(c) caption.]}}

%% ---------------- R1.8 ----------------
\subsection*{R1.8 --- ``Enhancement'' of $\Kzero = 1.0$ in the coordination-number plane (Fig.~12)}

\begin{rcomment}
In most cases, $N_L$ decreases monotonically with the increase of $\Kzero$,
as shown in Figures 8 and 10; however, the situation changes in Figure 12.
Theoretically, the evolution of the coordination number at the microscopic
level corresponds to the evolution of the macroscopic parameters such as the
excess pore water pressure ratio, deviatoric stress, and cumulative shear
work. The enhancement of $\Kzero = 1.0$ in the coordination number plane
should be explained more clearly.
\end{rcomment}

\response
\textit{[TODO: clarify physical cause --- isotropic consolidation distributes
contact normals more uniformly, yielding slightly higher $\Zmz$; reaffirm
that $\Zmz$ alone is insufficient to predict $N_L$ and that the directional
fabric indicator $\alpha$ is required, as demonstrated in
\secref{...}.]}

\changed{\textit{[TODO: 1--2 sentences appended to \figref{12} discussion.]}}

\clearpage

%% ===============================================================
%% Reviewer 2
%% ===============================================================

\section*{Response to Reviewer 2}

We thank Reviewer~2 for the thorough and constructive comments, and
particularly for pointing us to relevant prior work on DEM-HCA modelling.

%% ---------------- R2.1 ----------------
\subsection*{R2.1 --- Particle--wall friction coefficient set to zero (Table 1)}

\begin{rcomment}
Table 1: Why is the friction coefficient between particles and walls zero?
\end{rcomment}

\response
\textit{[TODO: explain that torque is transmitted geometrically through
blade faces rather than wall friction; setting wall friction to zero avoids
introducing an additional free parameter while the blade geometry already
ensures full torque transfer.]}

\changed{\textit{[TODO: Table~1 caption note / \secref{2.1}.]}}

%% ---------------- R2.2 ----------------
\subsection*{R2.2 --- Size of torsional blades}

\begin{rcomment}
What is the size of torsional blades?
\end{rcomment}

\response
\textit{[TODO: report blade dimensions from \texttt{generate.py}.]}

\changed{\textit{[TODO: \secref{2.1} or Table~1.]}}

%% ---------------- R2.3 ----------------
\subsection*{R2.3 --- Difference vs.\ servo mechanism in Han (2024)}

\begin{rcomment}
What is the difference between the servo mechanism used in this paper and
that in Han (2024)?
\end{rcomment}

\response
\textit{[TODO: reference the contribution statement in \secref{...}, lines
326--...; summarise: Han~(2024) solved only the radial and torsional
conditions with fixed specimen height; the present formulation adds the
axial degree of freedom and derives the cross-coupling
$\hat{\mathbf{K}}$ analytically.]}

\changed{\textit{[TODO: no new text needed, or minor clarification in
contribution paragraph.]}}

%% ---------------- R2.4 ----------------
\subsection*{R2.4 --- DEM-HCA literature review and the $H=D$ height choice}

\begin{rcomment}
There should be more context to introduce the research in DEM HCA models in
Part Introduction. \dots\ [boundary-condition types in Li~2014, Li~2017,
Song~2024a, Song~2024b; height $H \approx 2D$ in Shi~2021]. However, the
height equals to the outer diameter in this paper. Are there any
requirements or guidelines for the size of HCA?
\end{rcomment}

\response
\textit{[TODO: add Introduction paragraph reviewing prior DEM-HCA
modelling with a boundary-condition taxonomy (rigid wall, stacked wall,
flexible membrane); justify $H=D$ as a computational-cost trade-off with
PFC precedents; note that average fabric response, rather than end-plate
localisation, is the focus of this study.]}

\changed{\textit{[TODO: new paragraph in \secref{1}; short note in
\secref{2.1}.]}}

%% ---------------- R2.5 ----------------
\subsection*{R2.5 --- Stability of the servo near singular stiffness matrix}

\begin{rcomment}
The proposed servo mechanism may face numerical instability when the
stiffness matrix becomes near-singular due to the drastic loss of particle
contacts. How does the algorithm prevent numerical oscillations during
matrix inversion?
\end{rcomment}

\response
\textit{[TODO: depending on code audit --- either (a) describe
regularisation / condition-number checks, or (b) argue that the
constant-volume constraint eliminates radial-axis degeneracy, while the
torsional and axial degrees of freedom remain independently stable; any
post-liquefaction oscillations observed are physical rather than numerical.]}

\changed{\textit{[TODO: paragraph in \secref{...} or appendix.]}}

%% ---------------- R2.6 ----------------
\subsection*{R2.6 --- Evolution of $e_r$, $e_z$, or volume during validation}

\begin{rcomment}
For the validation part, please add some figures to display the evolution
of $e_r$ and $e_z$ over time or the change in volume.
\end{rcomment}

\response
\textit{[TODO: add plots of $r$, $H$, and total volume during monotonic
calibration and cyclic validation, confirming strict enforcement of the
constant-volume constraint.]}

\changed{\textit{[TODO: new panel in \figref{7} or a separate
\figref{...}.]}}

%% ---------------- R2.7 ----------------
\subsection*{R2.7 --- Contact-force anisotropy (3D rose diagrams)}

\begin{rcomment}
Fig.~14: You analysed contact density anisotropy, but why didn't you analyse
anisotropy for contact forces such as normal contact force? Please refer to
the papers mentioned before which plotted 3D rose diagrams of contact force
anisotropies.
\end{rcomment}

\response
\textit{[TODO: add contact-force rose diagrams alongside the contact-density
ones; reference Song~2024 precedent.]}

\changed{\textit{[TODO: extended \figref{14} or new companion figure.]}}

%% ---------------- R2.8 ----------------
\subsection*{R2.8 --- Quantifying the differences in Fig.~15}

\begin{rcomment}
Fig.~15: It is difficult to distinguish the difference in displacement
between $\Kzero=0.5$ and $2.0$ specimens when $N/N_L=0.00$ and $N/N_L=0.61$,
and it is also difficult to distinguish the difference in force chain when
$N/N_L=1.05$ and $N/N_L=1.07$. Please quantify the differences.
\end{rcomment}

\response
\textit{[TODO: add a quantitative metric --- e.g., displacement-magnitude
CDF and force-chain principal-direction distribution --- either as a new
panel or in the caption.]}

\changed{\textit{[TODO: \figref{15} caption / new panel.]}}

\clearpage

\section*{Closing}

We hope the revised manuscript now meets the journal's standards and the
reviewers' expectations. We remain grateful for the reviewers' careful
reading and valuable suggestions.

\bibliographystyle{plainnat}
\bibliography{../../refs.bib}

\end{document}
